% Supplementary item: hyperparameters_default
% GENERATED by scripts/generate_supp_table1.py --write. Do not edit;
% scripts/test_supp_table1.py fails on a hand edit. Rules: supplementary/README.md.

{\small
\begin{longtable}{@{}p{5.0cm}p{5.0cm}p{4.6cm}@{}}
\caption{\textbf{\SuppFile{} Table A: default hyperparameters for every model configuration in the study.}
Every value is read from \texttt{models/model\_defaults.py} at spec version 1.7.0, spec hash \texttt{9249926c38dc}, which is the file both the QM9 pipeline and the assay pipeline load their settings from. Every QM9 results row carries the spec version and the spec hash as two of its columns, so a QM9 table can be traced back to the settings that produced it; the three assay datasets are written by a second pipeline whose results files do not carry those two columns. The configuration list is the model roster the job generator submits. One setting is used for every representation and for every dataset, except where Table B gives a tuned replacement. The support vector machine uses an RBF kernel on every representation, so no result in this study carries a kernel that changes with the representation. Two values are not held in the shared spec and name their source file in the third column.}
\label{af1:hyperparameters}\\
\toprule
\textbf{Configuration} & \textbf{Hyperparameter} & \textbf{Default value} \\
\midrule
\endfirsthead
\toprule
\textbf{Configuration} & \textbf{Hyperparameter} & \textbf{Default value} \\
\midrule
\endhead
\multicolumn{3}{@{}l}{\textit{Shared by every neural network below}} \\
\midrule
Training, all networks & optimizer & adam \\
 & lr & 0.001 \\
 & weight\_decay & 0.0 \\
 & batch\_size & 32 \\
 & epochs & 100 \newline {\footnotesize\raggedright cap; early stopping usually reaches it first} \\
 & patience & 10 \newline {\footnotesize\raggedright epochs without a strictly better mean validation loss} \\
 & restore\_best\_weights & True \newline {\footnotesize\raggedright the reported fit is the best epoch, not the last} \\
 & val\_loss\_reduction & mean \\
 & improvement\_tolerance & 0.0 \\
 & mc\_passes & 100 \newline {\footnotesize\raggedright stochastic forward passes at inference} \\
 & standardise\_targets & True \newline {\footnotesize\raggedright fitted on training labels alone} \\
\midrule
\multicolumn{3}{@{}l}{\textit{Tree and boosting models}} \\
\midrule
RF \newline \footnotesize\texttt{rf} & n\_estimators & 100 \\
 & max\_depth & None \newline {\footnotesize\raggedright no limit} \\
 & min\_samples\_leaf & 5 \\
 & min\_samples\_split & 2 \\
 & max\_features & 0.3 \\
 & bootstrap & True \\
\addlinespace
RF300 \newline \footnotesize\texttt{rf300} & n\_estimators & 300 \\
 & max\_depth & None \newline {\footnotesize\raggedright no limit} \\
 & min\_samples\_leaf & 5 \\
 & min\_samples\_split & 2 \\
 & max\_features & 0.3 \\
 & bootstrap & True \\
\addlinespace
XGBoost \newline \footnotesize\texttt{xgboost} & n\_estimators & 100 \\
 & max\_depth & 6 \\
 & learning\_rate & 0.1 \\
 & subsample & 1.0 \\
 & colsample\_bytree & 1.0 \\
 & colsample\_bylevel & 1.0 \\
 & min\_child\_weight & 1 \\
 & gamma & 0.0 \\
 & reg\_alpha & 0.0 \\
 & reg\_lambda & 1.0 \\
\addlinespace
LightGBM \newline \footnotesize\texttt{lgb} & n\_estimators & 100 \\
 & learning\_rate & 0.1 \\
 & num\_leaves & 31 \\
 & max\_depth & $-1$ \newline {\footnotesize\raggedright no limit} \\
 & subsample & 1.0 \\
 & colsample\_bytree & 1.0 \\
 & min\_child\_samples & 20 \\
 & reg\_alpha & 0.0 \\
 & reg\_lambda & 0.0 \\
\addlinespace
NGBoost \newline \footnotesize\texttt{ngboost} & n\_estimators & 500 \\
 & learning\_rate & 0.01 \\
 & natural\_gradient & True \\
 & dist & Normal \newline {\footnotesize\raggedright ngboost.distns.Normal} \\
 & score & MLE \newline {\footnotesize\raggedright ngboost.scores.MLE, the log scoring rule} \\
 & minibatch\_frac & 1.0 \\
 & col\_sample & 1.0 \\
 & base\_max\_depth & 3 \newline {\footnotesize\raggedright depth of one boosting stage} \\
 & base\_criterion & friedman\_mse \\
 & early\_stopping\_rounds & 50 \newline {\footnotesize\raggedright stages without a better validation log likelihood} \\
 & use\_best\_iteration & True \newline {\footnotesize\raggedright predictions use the best stage, not every stage fitted} \\
 & ngboost\_ensemble\_seeds & 1 \newline {\footnotesize\raggedright 1 = one fit, no seed ensemble} \\
\addlinespace
QRF \newline \footnotesize\texttt{qrf} & n\_estimators & 300 \\
 & max\_depth & None \newline {\footnotesize\raggedright no limit} \\
 & min\_samples\_leaf & 5 \\
 & min\_samples\_split & 2 \\
 & max\_features & 0.3 \\
 & bootstrap & True \\
\midrule
\multicolumn{3}{@{}l}{\textit{Kernel models}} \\
\midrule
SVM \newline \footnotesize\texttt{svm} & C & 1.0 \\
 & gamma & scale \\
 & kernel & rbf \\
\addlinespace
GP \newline \footnotesize\texttt{gauche\_rbf} & kernel & rbf \newline {\footnotesize\raggedright set by this job, not by the shared spec} \\
 & likelihood\_noise & 0.001 \newline {\footnotesize\raggedright starting value; it is learned} \\
 & outputscale & 1.0 \\
 & apply\_outputscale & True \\
 & max\_train\_n & 5000 \newline {\footnotesize\raggedright molecules the process is fitted on; an exact Gaussian process is cubic in this. The test set is untouched and the number actually fitted is written into every results row} \\
 & fallback\_adam\_lr & 0.1 \newline {\footnotesize\raggedright used only when the standard fitter refuses the model} \\
 & fallback\_adam\_iters & 100 \newline {\footnotesize\raggedright used only when the standard fitter refuses the model} \\
 & single\_thread\_fit & False \newline {\footnotesize\raggedright a threading workaround, off since the environment check found one runtime} \\
 & init\_lengthscale\_from\_data & True \newline {\footnotesize\raggedright the RBF length scale starts at the median distance between training molecules, not at the library default} \\
 & lengthscale\_probe\_n & 500 \newline {\footnotesize\raggedright molecules sampled to estimate that median} \\
 & collapse\_fraction & 0.05 \newline {\footnotesize\raggedright a fit whose predictions vary by less than this fraction of the training label spread is recorded as collapsed rather than scored} \\
\addlinespace
GP-Tanimoto \newline \footnotesize\texttt{gauche} \newline \footnotesize run on ecfp4 only & kernel & tanimoto \newline {\footnotesize\raggedright set by this job, not by the shared spec} \\
 & likelihood\_noise & 0.001 \newline {\footnotesize\raggedright starting value; it is learned} \\
 & outputscale & 1.0 \\
 & apply\_outputscale & True \\
 & max\_train\_n & 5000 \newline {\footnotesize\raggedright molecules the process is fitted on; an exact Gaussian process is cubic in this. The test set is untouched and the number actually fitted is written into every results row} \\
 & fallback\_adam\_lr & 0.1 \newline {\footnotesize\raggedright used only when the standard fitter refuses the model} \\
 & fallback\_adam\_iters & 100 \newline {\footnotesize\raggedright used only when the standard fitter refuses the model} \\
 & single\_thread\_fit & False \newline {\footnotesize\raggedright a threading workaround, off since the environment check found one runtime} \\
 & init\_lengthscale\_from\_data & True \newline {\footnotesize\raggedright the RBF length scale starts at the median distance between training molecules, not at the library default} \\
 & lengthscale\_probe\_n & 500 \newline {\footnotesize\raggedright molecules sampled to estimate that median} \\
 & collapse\_fraction & 0.05 \newline {\footnotesize\raggedright a fit whose predictions vary by less than this fraction of the training label spread is recorded as collapsed rather than scored} \\
\addlinespace
GP-Hetero \newline \footnotesize\texttt{heteroscedastic\_gp} & kernel & rbf \newline {\footnotesize\raggedright set by this job, not by the shared spec} \\
 & likelihood\_noise & 0.001 \newline {\footnotesize\raggedright starting value; it is learned} \\
 & outputscale & 1.0 \\
 & apply\_outputscale & True \\
 & max\_train\_n & 5000 \newline {\footnotesize\raggedright molecules the process is fitted on; an exact Gaussian process is cubic in this. The test set is untouched and the number actually fitted is written into every results row} \\
 & fallback\_adam\_lr & 0.1 \newline {\footnotesize\raggedright used only when the standard fitter refuses the model} \\
 & fallback\_adam\_iters & 100 \newline {\footnotesize\raggedright used only when the standard fitter refuses the model} \\
 & single\_thread\_fit & False \newline {\footnotesize\raggedright a threading workaround, off since the environment check found one runtime} \\
 & init\_lengthscale\_from\_data & True \newline {\footnotesize\raggedright the RBF length scale starts at the median distance between training molecules, not at the library default} \\
 & lengthscale\_probe\_n & 500 \newline {\footnotesize\raggedright molecules sampled to estimate that median} \\
 & collapse\_fraction & 0.05 \newline {\footnotesize\raggedright a fit whose predictions vary by less than this fraction of the training label spread is recorded as collapsed rather than scored} \\
 & noise network hidden sizes & [64, 64] \newline {\footnotesize\raggedright models/models.py, not in the shared spec} \\
 & noise network output & Softplus, plus $10^{-4}$ \newline {\footnotesize\raggedright models/models.py, not in the shared spec} \\
 & Adam learning rate, process & 0.1 \newline {\footnotesize\raggedright models/models.py, not in the shared spec} \\
 & Adam learning rate, noise network & 0.001 \newline {\footnotesize\raggedright models/models.py, not in the shared spec} \\
 & noise network epochs & 100 \newline {\footnotesize\raggedright the shared epoch cap} \\
\midrule
\multicolumn{3}{@{}l}{\textit{Neural networks}} \\
\midrule
NN-$\alpha$ \newline \footnotesize\texttt{dnn} & hidden\_sizes & [128, 64] \\
 & dropout\_rate & 0.2 \\
 & activation & relu \\
\addlinespace
NN-$\beta$ \newline \footnotesize\texttt{mlp} & hidden\_size & 128 \\
 & num\_hidden\_layers & 2 \\
 & dropout\_rate & 0.2 \\
 & activation & relu \\
\addlinespace
BNN-$\alpha$ \newline \footnotesize\texttt{dnn\_bnn\_full} & hidden\_sizes & [128, 64] \\
 & dropout\_rate & 0.2 \\
 & activation & relu \\
 & bnn\_prior\_mu & 0.0 \\
 & bnn\_prior\_sigma & 0.1 \\
 & bnn\_kl\_weight & elbo \newline {\footnotesize\raggedright 1 / number of training molecules} \\
\addlinespace
BNN-$\beta$ \newline \footnotesize\texttt{mlp\_bnn\_full} & hidden\_size & 128 \\
 & num\_hidden\_layers & 2 \\
 & dropout\_rate & 0.2 \\
 & activation & relu \\
 & bnn\_prior\_mu & 0.0 \\
 & bnn\_prior\_sigma & 0.1 \\
 & bnn\_kl\_weight & elbo \newline {\footnotesize\raggedright 1 / number of training molecules} \\
\addlinespace
VBLL-$\alpha$ \newline \footnotesize\texttt{dnn\_bnn\_full\_variational} & hidden\_sizes & [128, 64] \\
 & dropout\_rate & 0.2 \\
 & activation & relu \\
 & vbll\_prior\_mu & 0.0 \\
 & vbll\_prior\_sigma & 1.0 \\
 & vbll\_init\_log\_sigma & $-3.0$ \\
 & vbll\_init\_log\_noise\_var & $-2.0$ \\
\addlinespace
VBLL-$\beta$ \newline \footnotesize\texttt{mlp\_bnn\_full\_variational} & hidden\_size & 128 \\
 & num\_hidden\_layers & 2 \\
 & dropout\_rate & 0.2 \\
 & activation & relu \\
 & vbll\_prior\_mu & 0.0 \\
 & vbll\_prior\_sigma & 1.0 \\
 & vbll\_init\_log\_sigma & $-3.0$ \\
 & vbll\_init\_log\_noise\_var & $-2.0$ \\
\addlinespace
VBLL-$\alpha$ (noise head) \newline \footnotesize\texttt{dnn\_bnn\_full\_variational\_hetero} & hidden\_sizes & [128, 64] \\
 & dropout\_rate & 0.2 \\
 & activation & relu \\
 & vbll\_prior\_mu & 0.0 \\
 & vbll\_prior\_sigma & 1.0 \\
 & vbll\_init\_log\_sigma & $-3.0$ \\
 & vbll\_init\_log\_noise\_var & $-2.0$ \\
 & observation noise & predicted per molecule \newline {\footnotesize\raggedright a variational noise head on the same features, started at the single-number value above} \\
\addlinespace
VBLL-$\beta$ (noise head) \newline \footnotesize\texttt{mlp\_bnn\_full\_variational\_hetero} & hidden\_size & 128 \\
 & num\_hidden\_layers & 2 \\
 & dropout\_rate & 0.2 \\
 & activation & relu \\
 & vbll\_prior\_mu & 0.0 \\
 & vbll\_prior\_sigma & 1.0 \\
 & vbll\_init\_log\_sigma & $-3.0$ \\
 & vbll\_init\_log\_noise\_var & $-2.0$ \\
 & observation noise & predicted per molecule \newline {\footnotesize\raggedright a variational noise head on the same features, started at the single-number value above} \\
\addlinespace
BNN-$\alpha$ (variance head) \newline \footnotesize\texttt{dnn\_bnn\_full\_mve} & hidden\_sizes & [128, 64] \\
 & dropout\_rate & 0.2 \\
 & activation & relu \\
 & bnn\_prior\_mu & 0.0 \\
 & bnn\_prior\_sigma & 0.1 \\
 & bnn\_kl\_weight & elbo \newline {\footnotesize\raggedright 1 / number of training molecules} \\
 & output head & value and log variance \newline {\footnotesize\raggedright two outputs} \\
 & loss & Gaussian negative log likelihood \\
 & log variance clamp & $[-10, 10]$ \newline {\footnotesize\raggedright scripts/loss\_functions.py, not in the shared spec} \\
\addlinespace
BNN-$\beta$ (variance head) \newline \footnotesize\texttt{mlp\_bnn\_full\_mve} & hidden\_size & 128 \\
 & num\_hidden\_layers & 2 \\
 & dropout\_rate & 0.2 \\
 & activation & relu \\
 & bnn\_prior\_mu & 0.0 \\
 & bnn\_prior\_sigma & 0.1 \\
 & bnn\_kl\_weight & elbo \newline {\footnotesize\raggedright 1 / number of training molecules} \\
 & output head & value and log variance \newline {\footnotesize\raggedright two outputs} \\
 & loss & Gaussian negative log likelihood \\
 & log variance clamp & $[-10, 10]$ \newline {\footnotesize\raggedright scripts/loss\_functions.py, not in the shared spec} \\
\bottomrule
\end{longtable}
}
